\documentclass[11pt,a4paper]{article}

\usepackage[T1]{fontenc}
\usepackage[utf8]{inputenc}
\usepackage{lmodern}
\usepackage[margin=2.5cm]{geometry}
\usepackage{amsmath}
\usepackage{booktabs}
\usepackage{longtable}
\usepackage{array}
\usepackage{xcolor}
\usepackage{fancyvrb}
\usepackage[colorlinks=true,linkcolor=blue!50!black,urlcolor=blue!50!black,citecolor=blue!50!black]{hyperref}

\usepackage{url}
% a keyword or file name (underscores allowed; long ones break after _ , / .)
\urlstyle{tt}
\makeatletter
\g@addto@macro{\UrlBreaks}{\do\_\do\,}
\makeatother
\DeclareUrlCommand\kw{}
% table of keywords: keyword, argument, default, description
\newenvironment{keywords}
 {\begin{longtable}{>{\raggedright}p{0.2\textwidth}>{\raggedright}p{0.12\textwidth}>{\raggedright}p{0.14\textwidth}p{0.41\textwidth}}
  \toprule keyword & argument & default & meaning \\ \midrule \endhead
  \bottomrule \endfoot}
 {\end{longtable}}
\newcommand{\key}[4]{\kw{#1} & #2 & #3 & #4 \\}

\title{\textbf{MonteGrappa}\\[2mm]\large Reference manual}
\author{G. Tiana, F. Villa, Y. Zhan, R. Capelli, C. Paissoni, P. Sormanni, R. Meloni}
\date{Version 1.3, 2026 revision}

\begin{document}
\maketitle

\begin{abstract}
MonteGrappa is a Monte Carlo program for the simulation of polymers and
proteins described by simplified, structure-based or statistical
potentials, from C$_\alpha$ models to all-atom models with rotamers. It
supports simulated and parallel tempering, periodic boxes with many chains,
and the iterative optimization of the potential against experimental
restraints. This manual describes the programs of the package, their input
files and every keyword they read. Please cite G.~Tiana \emph{et al.},
Comp.\ Phys.\ Comm.\ \textbf{186}, 93--104 (2015),
doi:\href{https://doi.org/10.1016/j.cpc.2014.09.012}{10.1016/j.cpc.2014.09.012}.
\end{abstract}

\tableofcontents

%======================================================================
\section{Overview}

The package contains five programs:
\begin{itemize}
\item \kw{montegrappa}, the Monte Carlo engine (serial, with simulated
  tempering, or with parallel tempering over MPI as \kw{montegrappa_mpi});
\item \kw{grappino}, which builds the input files of \kw{montegrappa} (polymer
  and potential) from a PDB structure;
\item \kw{mhistogram}, which computes thermodynamic quantities from energy
  histograms at several temperatures with the multiple-histogram method;
\item \kw{mgp2pdb}, which converts a polymer file into a PDB file;
\item \kw{dcd2pdb}, which converts a binary dcd trajectory into PDB frames.
\end{itemize}

A typical workflow is: prepare a PDB file, run \kw{grappino} to obtain the
polymer file (\kw{.pol}) and the potential file (\kw{.pot}), write a
parameter file (\kw{.par}) and run
\begin{Verbatim}
  montegrappa polymer.pol potential.pot parameters.par
\end{Verbatim}
The tutorials in \kw{tutorial/}, described in section~\ref{sec:tutorials},
go through these steps from a C$_\alpha$ Go model to parallel tempering,
simulated tempering and the optimization of the potential.

\paragraph{Units.} Lengths are in \AA{}ngstrom, angles in degrees.
Energies are in arbitrary units and the Boltzmann constant is 1, so
temperatures are in the same units as energies.

\paragraph{Input syntax.} All input files are plain text. In the parameter
files of \kw{montegrappa} and \kw{grappino} each line starts with a keyword
followed by its arguments; anything after them (e.g.\ a comment after
\kw{#}) is ignored. A keyword is recognized when the line \emph{starts}
with it, so keywords must begin at the first column; unknown keywords are
ignored silently, so check the log (each recognized keyword is echoed on
standard error). \emph{Flag} keywords switch an option on when they appear
alone and off when followed by \kw{0} (e.g.\ \kw{shell 0}). Polymer and
potential files are divided in chapters introduced by a line
\kw{[ name ]}.

%======================================================================
\section{Installation}

The code is written in C99 and needs only the math library; simulated
tempering and \kw{mhistogram} need GSL, parallel tempering needs MPI.
\begin{Verbatim}
  make                      serial montegrappa         -> bin/montegrappa
  make version=STEMPERING   with simulated tempering   -> bin/montegrappa
  make version=MPI          parallel tempering         -> bin/montegrappa_mpi
  make version=ASAN         AddressSanitizer + UBSan   -> bin/montegrappa_asan
  make grappino | mgp2pdb | dcd2pdb | mhistogram
  make all                  everything except ASAN
\end{Verbatim}
GSL is found with \kw{gsl-config}, MPI with the \kw{mpicc} wrapper; both can
be set on the command line, e.g.
\begin{Verbatim}
  make version=STEMPERING GSL_PREFIX=/opt/local
  make version=MPI MPICC=/path/to/mpicc
\end{Verbatim}
Each
version is compiled in its own directory \kw{obj/<version>}. The code is
always compiled with \kw{-ffp-contract=off}, so that the arithmetic (and
hence a trajectory with a given seed) is the same on all platforms.
\kw{make OPENMP=1} enables an experimental OpenMP parallelization of the
total energy.

\kw{make check} runs the test suite (about one minute): errors on wrong
input, bitwise regression of trajectories, consistency of the incremental
energy after every move, exact restart, cell list, tools, MPI, observables
and the distribution of the AVB move. \kw{make check-asan} runs it with the
sanitizers, \kw{make check-stat} tests detailed balance statistically (see
\kw{tests/README.md}).

%======================================================================
\section{Running montegrappa}

\begin{Verbatim}
  montegrappa  polymer.pol  potential.pot  parameters.par
  mpirun -np N  montegrappa_mpi  polymer.pol  potential.pot  parameters.par
  montegrappa  checkpoint  potential.pot  parameters.par      (restart)
\end{Verbatim}
With MPI, the number of processes must equal \kw{ntemp}: each rank
simulates one replica at one temperature.

\subsection{Monte Carlo steps and moves}
At each step, every move type whose counter has reached its frequency is
attempted once (see the move keywords in section~\ref{sec:moves}): a move
with frequency $N$ is attempted every $N$ steps. A move is accepted with the
Metropolis criterion $\min[1, e^{-\Delta E/T}]$, corrected by the ratio of
the proposal probabilities for the biased moves (\kw{movebias}, \kw{avbmc},
\kw{cbmc}). Moves that would break a constraint (locked dihedrals, bond
angles, box walls, overlaps) are rejected.

\subsection{Several runs}
\kw{nrun} $>1$ repeats the simulation; each run continues from the last
conformation of the previous one (from the initial one with
\kw{always_restart}). It is used together with the optimization of the
potential (section~\ref{sec:optim}), which changes the potential between
runs. Output files of run $r$ carry \kw{_run}$r$ in their name.

\subsection{Restart}
A checkpoint file (\kw{checkpoint.dat}, \kw{checkpoint_proc<rank>.dat}
with MPI) is written every \kw{tcheckpoint} seconds of wall-clock time
(default 600) or every \kw{ncheckpoint} steps; with MPI every 100000 steps
by default, so that all replicas write it at the same step. Each run also
writes one when it starts. It contains the state of the random number
generator, all counters and the coordinates with 17 significant digits, and
the length of the output files at that moment. The command
\begin{Verbatim}
  montegrappa checkpoint potential.pot parameters.par
\end{Verbatim}
in the same directory continues the simulation exactly (bit by bit) as if
it had not been interrupted: the trajectory, energy, observables and
exchange files are cut back to their length at the checkpoint, so nothing
is written twice, and the log is continued. The original polymer file must
still be there. \kw{kill -USR1} stops the program cleanly at the end of the
current step after writing a checkpoint. The state of simulated tempering
and the frames recorded for the optimization of the potential are not
saved in the checkpoint.

\subsection{Output files}
\begin{longtable}{>{\raggedright}p{0.32\textwidth}p{0.6\textwidth}}
\toprule file & content \\ \midrule
\kw{traj.pdb} (keyword \kw{traj}) & trajectory in PDB format every \kw{nprinttrj} steps, with the initial and final conformations; with \kw{trjformat dcd} a binary \kw{traj.dcd} and its topology \kw{traj_top.pdb} \\
\kw{energy.ene} (\kw{efile}) & every \kw{nprinte} steps: step, total energy, pair energy, angle energy, dihedral energy, H-field energy \\
\kw{last.pol} (\kw{lastp}) & last conformation, as a polymer file \\
\kw{montegrappa.log} (\kw{logfile}) & log of the input and of the run \\
\kw{obs.dat} & aggregation observables every \kw{nprintobs} steps \\
\kw{mginfo} & progress of the simulation (as Amber's \kw{mdinfo}) every \kw{tmginfo} seconds \\
\kw{checkpoint.dat} & checkpoint for the restart \\
\kw{proc_<rank>.log} (\kw{procfile}) & with MPI, log of each rank \\
\kw{exchange_proc<rank>.dat} & with MPI, accepted exchanges (\kw{[EX] step i i+1}) \\
\kw{newpot_<run>.pot}, \kw{chi2.dat} & optimized potential of each run and $\chi^2$ (section~\ref{sec:optim}) \\
\bottomrule
\end{longtable}
With MPI the names carry \kw{_proc<rank>}, with \kw{nrun}$>1$
\kw{_run<run>}. The PDB trajectory has one \kw{TITLE} line and one
\kw{ENDMDL} per frame; chains are named A, B, \dots{} (and separated by
\kw{TER} from the 25th chain on). Coordinates are not wrapped into a
periodic box; far from the origin they keep their 8 columns with fewer
decimals.

%======================================================================
\section{The polymer file}
\label{sec:pol}

The polymer file describes the chains: the backbone atoms with their
coordinates, the rotamers of the sidechains and the current rotameric
state. It is normally written by \kw{grappino}.

\paragraph{\texttt{[ backbone ]}} One line per backbone atom:
\begin{Verbatim}
  back  ia  type  itype  aa  iaa  ch  x  y  z  tomove
\end{Verbatim}
\kw{back} is the index of the atom along the backbone of its chain (from
0), \kw{ia} a global atom identifier (unique over the whole system),
\kw{type} the atom name, \kw{itype} the atom type used by the potential,
\kw{aa} and \kw{iaa} the residue name and number, \kw{ch} the chain (from
0), \kw{x y z} the coordinates and \kw{tomove} 1 if the dihedral around
this atom can change, 0 if it is locked (e.g.\ the peptide bond).

\paragraph{\texttt{[ rotamers ]}} The sidechain atoms attached to each backbone
atom, for each rotamer:
\begin{Verbatim}
  back  ch  rot  at   b1  b2  b3   ia  type  itype   ang  dih  r
\end{Verbatim}
Sidechain atom \kw{at} of backbone atom \kw{back} of chain \kw{ch}, in
rotamer \kw{rot}, is placed at distance \kw{r}, angle \kw{ang} and
dihedral \kw{dih} from the atoms with identifiers \kw{b1 b2 b3} (internal
coordinates). \kw{ia}, \kw{type} and \kw{itype} are as in the backbone.

\paragraph{\texttt{[ sidechains ]}} The current rotamer of each residue:
\kw{back ch irot}. A residue whose rotamer is not listed is in rotamer 0;
the rotamer must exist.

%======================================================================
\section{The potential file}
\label{sec:pot}

The energy is the sum of pair interactions between atoms of different
residues, and of angular, dihedral and H-field terms:
\begin{multline*}
E = \sum_{\text{pairs}} \epsilon_{ab}\,\theta(r_{w,ab} - r)
 + \sum_i k_i(\vartheta_i-\vartheta_i^0)^2 \\
 + \sum_i \left[ k^{(1)}_i\left(1-\cos(\phi_i-\phi^{(1)}_i)\right)
               + k^{(3)}_i\left(1-\cos 3(\phi_i-\phi^{(3)}_i)\right)\right]
 + E_{\text{HF}}
\end{multline*}
A pair of atoms of types $a,b$ at distance $r$ has energy
$\epsilon_{ab}$ if $r<r_{w,ab}$ (square well), and infinite energy if
$r<r_{0,ab}$ (hard core). Atoms of residues closer than \kw{imin} along
the same chain do not interact, except through the global hard core.

\paragraph{\texttt{[ pairs ]}} \texttt{i j e rw r0}: atom types $i$, $j$, well
depth $\epsilon$, well radius $r_w$, hard-core radius $r_0$. The matrix is
symmetric. Types that no atom of the polymer has are ignored, so a general
potential can be used for any polymer.

\paragraph{\texttt{[ angles ]}} \texttt{ia k ang0}: harmonic bond-angle term
centred on the backbone atom with identifier \kw{ia}.

\paragraph{\texttt{[ dihedrals ]}} \texttt{ia k1 phi01 k3 phi03}: periodic dihedral
term (see above) on the dihedral around backbone atom \kw{ia}.

\paragraph{\texttt{[ global ]}} Keywords applying to all atoms:
\begin{keywords}
\key{hardcore}{r0}{0}{hard-core radius of all pairs (also those closer than \kw{imin})}
\key{homopolymeric}{e r}{--}{square well of depth $e$ and radius $r$ between all types}
\key{imin}{n}{0}{atoms whose backbone indices along the same chain differ by $n$ or less do not interact (only the global hard core)}
\key{angle}{k ang0}{--}{the same angular term for all atoms}
\key{dihedral1}{k phi0}{--}{the same $n=1$ dihedral term for all atoms}
\key{dihedral3}{k phi0}{--}{the same $n=3$ dihedral term for all atoms}
\key{splice}{kr ke}{off}{beyond $kr\cdot r_w$ ($kr<1$) the well has depth $ke\cdot\epsilon$ (two-step well); with hydrogen bonds, donor--acceptor pairs have no outer step}
\key{e_dih}{e}{--}{strength of the tabulated dihedral potential}
\key{g_dihbin}{bin}{--}{bin (degrees) of the tabulated dihedrals}
\key{e_dihram}{e}{1}{strength of the Ramachandran dihedral potential}
\key{boxtype}{c, s, p}{none}{cubic box with walls, spherical box, periodic cubic box}
\key{boxsize}{L}{--}{side of the cubic box, diameter of the spherical one (\AA)}
\end{keywords}
In a box with walls (\kw{c}: cube of side $L$ centred in the origin;
\kw{s}: sphere of diameter $L$) backbone atoms cannot leave the box. In a
periodic box (\kw{p}) pairs interact through the minimum image; $L$ must
exceed twice the range of interaction of two residues, and shells
(\kw{shell}) cannot be used.

\paragraph{\texttt{[ hydrogen\_bonds ]}} \texttt{ia d|a ib}: atom \texttt{ia} is a donor
(\kw{d}) or an acceptor (\kw{a}) and \kw{ib} is the atom bonded to it that
defines its direction. With \kw{hb} in the parameter file, the well of a
donor--acceptor pair is multiplied by
$\sqrt{|\cos\vartheta_1\cos\vartheta_2|}$, where $\vartheta_{1,2}$ are the
angles between the bonds and the donor--acceptor vector.

\paragraph{\texttt{[ dihedral\_weight ]}, \texttt{[ dihedral\_table ]}} Tabulated
dihedral potential: \kw{ia pa pb which} assigns to atom \kw{ia} the weights
of the $\alpha$ and $\beta$ tables (\kw{which} 0 for $\phi$, 1 for $\psi$);
the table has lines \kw{bin phi_a phi_b psi_a psi_b}.

\paragraph{\texttt{[ Ramachandran\_Dihedrals ]}, \texttt{[ Alfa/Beta\_propensity ]}}
Gaussian Ramachandran potential: lines \kw{a|b f|p sigma dih0} give width
and centre for $\alpha$/$\beta$ and $\phi$/$\psi$; lines \kw{iaa pa pb} the
propensities of each amino acid.

\paragraph{\texttt{[ hfields ]}} \texttt{itype h}: H-field energy $h$ of a residue
whose first sidechain atom has type \kw{itype}, multiplied by its number of
contacts. A potential without this chapter switches H fields off.

\paragraph{\texttt{[ comments ]}} Free text, printed on standard error.

%======================================================================
\section{The parameter file of montegrappa}
\label{sec:par}

\subsection{System and run}
\begin{keywords}
\key{nchains}{n}{required}{number of chains in the polymer file}
\key{nstep}{n}{100000}{Monte Carlo steps of each run}
\key{Temp}{T}{required}{temperature (serial and simulated tempering)}
\key{seed}{n}{clock}{seed of the random number generator (xoshiro256**)}
\key{nrun}{n}{1}{number of runs (iterations of the optimization)}
\key{always_restart}{flag}{off}{every run starts from the initial conformation}
\key{nosidechain}{flag}{off}{ignore the sidechains}
\key{noangpot}{flag}{off}{no angular potential}
\key{nodihpot}{flag}{off}{no dihedral potential}
\key{nohfields}{flag}{off}{no H fields}
\key{r_contact}{r}{5}{contact distance (\AA) for the H fields}
\key{hb}{flag}{off}{hydrogen bonds (chapter \kw{[ hydrogen_bonds ]} of the potential)}
\key{disentangle}{flag}{off}{overlapping atoms are allowed with an energy of $10^4$ per pair (kept as a pair energy), so that a structure with overlaps can relax; for preparing structures, not for production}
\key{shell}{flag}{off}{use a neighbour list of radius \kw{r2shell}}
\key{r2shell}{r}{6}{radius (\AA) of the neighbour list (increased if shorter than the interactions)}
\key{nshell}{n}{10}{steps between updates of the neighbour list}
\key{cells}{1, 0, -1}{-1}{cell list for the energy: always, never, automatic (for large systems)}
\end{keywords}

\subsection{Moves}
\label{sec:moves}
The value of each move keyword is its frequency: the move is attempted
every $n$ steps. Moves not listed are not used; at least one is required.
\begin{keywords}
\key{flip}{n}{--}{rotate a backbone atom around the axis of its two neighbours, by up to $\pm$\kw{dw_flip}}
\key{pivot}{n}{--}{rotate one end of a chain around a random backbone bond, by up to $\pm$\kw{dw_pivot}/2}
\key{mpivot}{n}{--}{\kw{nmul_mpivot} consecutive pivots, by up to $\pm$\kw{dw_mpivot}/2}
\key{lpivot}{n}{--}{local (loose) pivot: \kw{nmul_lpivot} consecutive pivots by up to $\pm$\kw{dw_lpivot}/2 with the rest of the chain fixed; the squared length of the bond that closes the moved segment may change by up to \kw{r_cloose}$^2$}
\key{mflip}{n}{--}{flip of a piece of chain up to \kw{nmul_mflip} atoms long, by up to $\pm$\kw{dw_mflip}}
\key{sidechain}{n}{--}{change the rotamer of a residue}
\key{movebias}{n}{--}{biased Gaussian local move (Favrin \emph{et al.}) on \kw{nmul_local} atoms, with \kw{bgs_a}, \kw{bgs_b}}
\key{backrub}{n}{--}{rotate a segment of five backbone atoms around the axis of the C$_\alpha$ at its ends}
\key{backsiderub}{n}{--}{backrub with a new rotamer of the central residue}
\key{movecom}{n}{--}{translate a chain by up to $\pm$\kw{dx_com}/2 per axis}
\key{moverot}{n}{--}{rotate a chain by up to $\pm$\kw{dtheta}/2}
\key{comcluster}{n}{--}{translate a cluster of chains in contact by up to $\pm$\kw{dx_clm}/2 per axis}
\key{rotcluster}{n}{--}{rotate a cluster of chains by up to $\pm$\kw{dtheta}/2}
\key{avbmc}{n}{--}{aggregation-volume-bias move (Chen and Siepmann, J.\ Phys.\ Chem.\ B \textbf{105}, 11275 (2001)): a chain is moved, with a random orientation, into or out of the region within \kw{avb_rin} of another chain; periodic box only}
\key{cbmc}{n}{--}{configurational-bias regrowth of the last 1..\kw{nmax_cbmc} backbone atoms of a chain end with \kw{ntrial_cbmc} trials per dihedral (Siepmann and Frenkel, Mol.\ Phys.\ \textbf{75}, 59 (1992))}
\end{keywords}

\noindent Parameters of the moves:
\begin{keywords}
\key{dw_flip}{deg}{30}{amplitude of \kw{flip}}
\key{dw_pivot}{deg}{10}{amplitude of \kw{pivot}}
\key{dw_mpivot}{deg}{1}{amplitude of \kw{mpivot}}
\key{dw_lpivot}{deg}{1}{amplitude of \kw{lpivot}}
\key{dw_mflip}{deg}{30}{amplitude of \kw{mflip}}
\key{randdw}{1, 2}{1}{pivot angles uniform (1) or Gaussian of that width (2)}
\key{nmul_mpivot}{n}{3}{pivots of \kw{mpivot}}
\key{nmul_lpivot}{n}{3}{pivots of \kw{lpivot} (negative: random between 2 and $|n|$)}
\key{nmul_mflip}{n}{100}{largest length of \kw{mflip}}
\key{nmul_local}{n}{required}{atoms moved by \kw{movebias} (at least 3)}
\key{bgs_a, bgs_b}{a b}{2, 1}{parameters of the Gaussian of \kw{movebias}}
\key{dx_com}{\AA}{1}{amplitude of \kw{movecom}}
\key{dx_clm}{\AA}{1}{amplitude of \kw{comcluster}}
\key{dtheta}{deg}{1}{amplitude of \kw{moverot} and \kw{rotcluster}}
\key{avb_rin}{\AA}{required}{radius of the bonded region of \kw{avbmc} (box side $>2\cdot$\kw{avb_rin})}
\key{avb_pin}{p}{0.5}{probability of a move into the bonded region}
\key{nmax_cbmc}{n}{6}{largest number of atoms regrown by \kw{cbmc}}
\key{ntrial_cbmc}{n}{8}{trials per dihedral of \kw{cbmc}}
\key{r_cloose}{\AA}{0.5}{tolerance on the closing bond of \kw{lpivot}: $|d^2-d_0^2|<$\kw{r_cloose}$^2$}
\key{a_cloose}{deg}{off}{largest change of a bond angle allowed to \kw{flip}, \kw{lpivot}}
\key{d_cloose}{deg}{off}{largest change of a locked dihedral (\kw{tomove 0}) allowed to flips, \kw{lpivot}, \kw{mflip}, \kw{backrub}}
\end{keywords}

\subsection{Output}
\begin{keywords}
\key{nprinttrj}{n}{1000}{steps between frames of the trajectory}
\key{nprinte}{n}{1000}{steps between lines of the energy file}
\key{nprintlog}{n}{1000}{steps between lines of the log}
\key{nprintobs}{n}{0}{steps between aggregation observables (0: none)}
\key{traj}{name}{traj}{base name of the trajectory}
\key{trjformat}{pdb, dcd}{pdb}{format of the trajectory}
\key{efile}{name}{energy}{base name of the energy file}
\key{lastp}{name}{last}{base name of the last conformation}
\key{logfile}{name}{\kw{montegrappa.log}}{log file}
\key{procfile}{name}{proc}{base name of the logs of the MPI ranks}
\key{ncheckpoint}{n}{0}{steps between checkpoints (0: by time)}
\key{tcheckpoint}{s}{600}{seconds between checkpoints}
\key{tmginfo}{s}{60}{seconds between updates of \kw{mginfo} (0: not written)}
\key{echeck}{flag}{off}{after every move compare the energy, contacts, cell list and sidechains with a full recalculation (slow; to test moves)}
\key{debug}{n}{0}{verbosity (needs compilation with \kw{-DDEBUG})}
\end{keywords}
The aggregation observables are: number of clusters of chains in contact,
largest and mean cluster size, pairs of chains in contact, parallel and
antiparallel pairs (cosine of the end-to-end vectors above 0.8), pairs in
register and mean register, nematic order parameter.

\subsection{Parallel tempering (montegrappa\_mpi)}
\begin{keywords}
\key{ntemp}{n}{required}{number of replicas (= MPI processes)}
\key{temperatures}{}{required}{followed by \kw{ntemp} lines, the temperature of each rank}
\key{step_exchange}{n}{10000}{steps between exchanges (at least 2): even pairs at half the interval, odd pairs at its end}
\key{models}{n}{1}{1: all replicas start from the polymer file; \kw{ntemp}: each from its own file}
\key{replicas}{}{--}{followed by \kw{ntemp} lines, the polymer file of each rank (with \kw{models ntemp})}
\key{iT_bias}{n}{--}{with shells, ranks from $n$ on do not use \kw{movebias}}
\end{keywords}
Exchanges between ranks $i$ and $i+1$ are accepted with
$\min\{1,\exp[(1/T_i-1/T_{i+1})(E_i-E_{i+1})]\}$. Each rank has its own
independent stream of random numbers.

\subsection{Simulated tempering (version STEMPERING)}
Simulated tempering is switched on with \kw{stempering}. The temperature
changes among a ladder with weights $g$; with \kw{st_method adaptive} the
ladder and weights are adjusted during the run from the density of states
computed with the multiple-histogram method.
\begin{keywords}
\key{stempering}{flag}{off}{switch on simulated tempering}
\key{st_method}{stempering, adaptive}{--}{fixed or adaptive ladder}
\key{st_ntemp}{n}{1}{number of temperatures}
\key{st_temperatures}{}{--}{followed by \kw{st_ntemp} lines \kw{T [g]}}
\key{st_ntmax}{n}{30}{largest number of temperatures}
\key{st_nstep}{n}{--}{steps between attempts to change temperature}
\key{st_nsadj}{n}{100000}{steps between adjustments of the ladder}
\key{st_preamble}{n}{0}{steps before the histograms are collected}
\key{st_emin, st_emax, st_ebin}{E}{--}{range and bin of the energy histograms}
\key{st_keepall}{flag}{off}{use all past histograms in the multiple-histogram method}
\key{st_nkeepold, st_sumoldhisto}{n}{0}{how many old histograms to keep and whether to sum them}
\key{st_hthresh, st_binthresh, st_phthresh}{x}{0.7, 0, 0.1}{thresholds on histogram overlap, bin counts and probabilities}
\key{st_k_new, st_lp_new, st_anneal}{}{1, 1, --}{how new temperatures are added (\kw{st_anneal sigma} or \kw{prob})}
\key{st_lpthresh, st_t_only_add, st_removet}{}{0, 0, 99}{when temperatures are added or removed}
\key{st_tstop, st_ttarget, st_tnorm}{T}{0}{stop / target / normalization temperatures}
\key{st_ignoreb, st_deltat, st_gmethod}{}{0, 0, 0}{as \kw{ignoreb} and \kw{deltat} of \kw{mhistogram}; method for the weights}
\key{st_nfail1, st_nfail2}{n}{0, 20}{tolerated failures of the adjustment}
\key{st_paranoid}{flag}{off}{stop on warnings}
\key{st_restart}{n}{-1}{write a restart of the ladder}
\key{st_nprint, st_nprntt, st_printpdb}{n}{1, 500000, 1000}{output intervals}
\key{st_tfile}{name}{\kw{temperatures.dat}}{temperatures visited}
\key{st_thefile}{name}{\kw{thermdodyn.dat}}{thermodynamics}
\key{st_dosfile}{name}{\kw{dos.dat}}{density of states}
\key{st_histofile}{name}{--}{histograms}
\key{st_dumbfile}{name}{\kw{dumb.dat}}{average energies at each temperature}
\key{st_pdbnf}{name}{\kw{harvest.pdb}}{structures harvested at \kw{st_ttarget}}
\key{st_debug}{n}{0}{verbosity}
\end{keywords}

\subsection{Optimization of the potential}
\label{sec:optim}
With \kw{op_minim sample}, at the end of each run the matrix of the pair
energies $\epsilon_{ab}$ is optimized so that the averages of the
restraints in \kw{op_file}, reweighted from the frames recorded during the
run, match their experimental values; the new potential is written to
\kw{newpot_<run>.pot} and used in the next run (\kw{nrun}). The frames are
reweighted to temperature \kw{op_T}; a step changes one element by a
random amount up to $\pm$\kw{op_step}/2 within
[\kw{op_emin},\kw{op_emax}] and is accepted if it lowers $\chi^2$ (or with
the Metropolis criterion at \kw{op_tchi2}).
\begin{keywords}
\key{op_minim}{none, sample}{none}{switch on the optimization}
\key{op_file}{name}{--}{restraints (see below)}
\key{op_T}{T}{1}{temperature of the reweighting}
\key{op_deltat}{n}{10000}{steps between recorded frames}
\key{op_wait}{n}{0}{steps before the first frame}
\key{op_itermax}{n}{1000}{steps of the optimization}
\key{op_step}{e}{1}{largest change of an element}
\key{op_stop}{$\chi^2$}{0.1}{stop when $\chi^2$ is lower}
\key{op_emin, op_emax}{e}{$\mp9\cdot10^{19}$}{bounds of the elements}
\key{op_tchi2}{T}{0}{fictitious temperature of the search on the reduced $\chi^2$ (0: only steps lowering it); the best matrix is returned}
\key{op_neff}{F}{0}{smallest effective fraction of frames allowed in the reweighting, relative to its initial value}
\key{op_allreplicas}{flag}{off}{with MPI, all replicas record frames and rank 0 reweights them with MBAR}
\key{op_rw, op_r0}{\AA}{3, 1}{well and hard-core radius of pairs not defined in the potential}
\key{op_print}{n}{100}{steps between lines of the log}
\key{record_native}{flag}{off}{the starting structure is the first frame}
\end{keywords}
The \kw{op_file} starts with \kw{ndata n} followed by $n$ restraints
\begin{Verbatim}
  i1 i2 type value sigma            or   i1 i2 i3 i4 type value sigma
\end{Verbatim}
with \kw{type} 0: contact between residues \kw{i1} and \kw{i2}; 1:
distance between atoms \kw{i1} and \kw{i2}; 2: $1/r^6$ between them; 3:
contact between the blocks of residues \kw{i1..i2} and \kw{i3..i4}.
Restraints with negative \kw{sigma} are not fitted: their $\chi^2$ is a
validation, written in the third column of \kw{chi2.dat}.

%======================================================================
\section{grappino}
\label{sec:grappino}

\begin{Verbatim}
  grappino input.in
\end{Verbatim}
reads a PDB file and writes the polymer and potential files (and
optionally a restraint file for the optimization). The model is defined by
the backbone atoms kept, the sidechains (from the PDB, from a rotamer
library, or none), the atom types and the potential.

\begin{keywords}
\key{pdbfile}{name}{required}{input structure}
\key{polfile}{name}{\kw{polymer.pol}}{output polymer file}
\key{potfile}{name}{\kw{potential.pot}}{output potential file}
\key{model}{CA, CACB, NCAC, BackCB}{all}{simplify the PDB to these atoms}
\key{backbone_atoms}{n names}{3 N CA C}{backbone atoms (e.g.\ \kw{1 CA})}
\key{locked_atoms}{names}{--}{backbone atoms whose dihedral is locked (after \kw{backbone_atoms})}
\key{hydrogens}{flag}{off}{keep the hydrogens of the PDB}
\key{nosidechain}{flag}{off}{no sidechains}
\key{rotamers}{flag}{off}{sidechains from the rotamer library \kw{rotfile}}
\key{rotfile}{name}{--}{rotamer library (e.g.\ \kw{lib/rotamers.lib})}
\key{pdb_rot}{flag}{off}{add the rotamer of the PDB to the library ones}
\key{cb_pdb}{flag}{on}{keep the C$_\beta$ of the PDB (0: pseudo-C$_\beta$ in the backbone plane)}
\key{sidescratch}{flag}{off}{build all sidechains from the library}
\key{atomtypes}{\kw{go}, \kw{go_aa}, \kw{go_homo}, file}{required}{\kw{go}: one type per atom; \kw{go_aa}: one type per residue; \kw{go_homo}: types shared by identical chains (e.g.\ fibrils); a file: types from a library}
\key{potential}{go, coevo, pairs}{go}{native-centric Go potential, from coevolution data (\kw{coevofile}), or from a matrix of pairs (\kw{pairsfile})}
\key{go-energy}{e}{-1}{well depth of native contacts}
\key{r_native}{\AA}{3.5}{atoms closer than this in the PDB are in native contact}
\key{r_hardcore}{\AA}{2}{hard-core radius}
\key{imin}{n}{9}{backbone atoms of the same chain whose indices differ by $n$ or less are not in contact (as in montegrappa)}
\key{r_bonded}{\AA}{2}{threshold for bonded atoms}
\key{use_nativedist}{flag}{off}{well radius from the native distance: \kw{k_native_r} times it, hard core \kw{k_native_hc} times it}
\key{k_native_r, k_native_hc}{x}{1, 0.5}{factors of \kw{use_nativedist}}
\key{e_homo, r_homo}{e r}{0, 0}{global homopolymeric well}
\key{splice}{kr ke}{off}{two-step well (see \kw{[ global ]})}
\key{go_dihedrals}{flag}{off}{dihedral terms centred on the native dihedrals with \kw{e_dih1}, \kw{e_dih3}}
\key{go_angles}{flag}{off}{angular terms centred on the native angles with \kw{e_ang}}
\key{dihfile, e_dihtab}{name e}{--, 1}{tabulated statistical dihedrals}
\key{dih_ram}{flag}{off}{Ramachandran dihedrals (\kw{e_dihram}, \kw{sig_*}, \kw{phi_0_*}, \kw{psi_0_*}, \kw{dih_ka}, \kw{dih_kb}, \kw{propensityfile})}
\key{hb_backbone}{flag}{off}{backbone hydrogen bonds, with amide hydrogens built at \kw{hb_dist}}
\key{h_fields}{flag}{off}{H fields from \kw{h_fieldsfile} and \kw{maxcontfile}, scaled by \kw{hfs_alpha}}
\key{cys_e, cys_r}{e r}{0}{disulphide bridges}
\key{go_noise, go_nsigma}{n s}{0, 1}{random noise on the Go energies}
\key{pairsfile}{name}{--}{matrix of pairs of \kw{potential pairs}}
\key{coevofile}{name}{--}{coevolution data of \kw{potential coevo} (\kw{coevo_allatom}, \kw{contactfile}, \kw{maxcontpairsfile})}
\key{op_file, op_kind}{name kind}{--, \kw{GO_DIST_CA}}{write a restraint file (native C$_\alpha$ distances, or all atoms with \kw{GO_DIST_ALLATOM})}
\key{debug}{n}{0}{verbosity}
\end{keywords}

%======================================================================
\section{Other tools}

\paragraph{mhistogram.} \texttt{mhistogram input} computes the density of
states and the thermodynamics (average energy, specific heat, free energy)
from energy histograms (or energy series with \kw{nohisto}) at several
temperatures, solving the multiple-histogram equations (G.~Tiana and
L.~Sutto, Phys.\ Rev.\ E \textbf{84}, 061910 (2011)). The input lists
\kw{ntemp}, then after \kw{files} one line per temperature
(\kw{file T column_energy column_counts}); options include \kw{ebin},
\kw{emin}, \kw{threshold}, \kw{outfile}, \kw{gfile}, \kw{zfile},
\kw{tmin}, \kw{tbin}, \kw{ntbin}, \kw{kb}, \kw{ignoreb}, \kw{deltat}.
\kw{mhistogram -h} lists them all.

\paragraph{mgp2pdb.} \texttt{mgp2pdb -i polymer.pol -o polymer.pdb} converts a
polymer file to PDB.

\paragraph{dcd2pdb.} \kw{dcd2pdb traj_top.pdb traj.dcd [first [last
[stride]]] > frames.pdb} converts a dcd trajectory to PDB models; the
topology file is written by \kw{montegrappa} with the dcd file.

%======================================================================
\section{Tutorials}
\label{sec:tutorials}

The directory \kw{tutorial/} contains six worked examples, each with its
input files, a \kw{README} with the commands to run, and (except 5 and 6)
the results obtained with this version of the code in \kw{results/}. Run
the commands from inside the directory of the tutorial, with the paths of
the executables in \kw{bin/}. A run with the same seed (and, with MPI, the
same number of processes) reproduces them exactly, except tutorial 2, whose
parameter file takes the seed from the clock; the RMSD were computed
with \kw{tools/rmsd.py} (C$_\alpha$ atoms after optimal superposition, in
nm, as \kw{gmx rms}).

\subsection{Tutorial 1: unfolding and refolding of a C\texorpdfstring{$_\alpha$}{-alpha} Go model}
\emph{Programs:} \kw{grappino}, \kw{montegrappa}. \emph{System:} the 64
residues of chymotrypsin inhibitor 2 (\kw{1YPC.pdb}).

This is the basic workflow. \kw{1YPC-CA.in} tells \kw{grappino} to keep only
the C$_\alpha$ atoms (\kw{backbone_atoms 1 CA}, \kw{nosidechain}), to give
each atom its own type (\kw{atomtypes go}) and to build a Go potential: every
pair of C$_\alpha$ closer than \kw{r_native} = 6~\AA{} in the native
structure and more than \kw{imin} = 4 residues apart along the chain gets a
well of depth $-1.5$, all pairs a hard core of 4~\AA; \kw{go_dihedrals}
adds dihedral terms centred on the native dihedrals. The command
\begin{Verbatim}
  grappino 1YPC-CA.in
\end{Verbatim}
writes \kw{1YPC-CA.pol} and \kw{1YPC-CA.pot}. The first simulation,
\begin{Verbatim}
  montegrappa 1YPC-CA.pol 1YPC-CA.pot 1YPC-CA.unf.par
\end{Verbatim}
runs $5\cdot10^6$ steps at $T=1.6$, above the folding temperature, with
pivot (every 4 steps) and multiple-pivot (every step) moves: the energy in
\kw{energy.ene} rises and the RMSD from the native structure grows as the
protein unfolds. The second one starts from the last, unfolded,
conformation and runs at $T=1$, where the native state is stable:
\begin{Verbatim}
  montegrappa last.pol 1YPC-CA.pot 1YPC-CA.fold.par
\end{Verbatim}
The energy drops and the RMSD decreases to a few \AA{}ngstrom. To compute
the RMSD, use the first frame of \kw{traj.pdb} as reference rather than
\kw{1YPC.pdb}, because the trajectory contains only the atoms of the model,
numbered in its own order. \kw{results/} contains the energies and the RMSD
of both runs. With seed 42 the chain refolds only partially within $5\cdot10^6$
steps (RMSD about 2~\AA): refolding from an unfolded state is a rare event
on this time scale, and other seeds refold completely or not at all.

\subsection{Tutorial 2: a chromatin polymer from 5C/Hi-C contact data}
\emph{Programs:} \kw{montegrappa} (potential optimization).
\emph{System:} a bead-spring polymer of 20 beads.

\kw{expdata.dat} contains contact counts between pairs of beads
(\kw{bead1 bead2 average stdev}), as measured by a 5C or Hi-C experiment.
The script \kw{generate_5C.sh} asks for the data file, the number of beads,
a normalization constant, a root name, the energy scale of the initial
interaction matrix, the interaction range and the hard-core radius (the
\kw{README} gives typical answers), and writes:
\begin{itemize}
\item \kw{test.pol}, a chain of beads one unit apart, each bead of its own type;
\item \kw{test.op}, the restraints: for each pair, a contact (type 0) whose
  target probability $p$ is the count divided by the normalization constant
  (at most 1), with the standard deviation divided likewise (at least 0.01)
  as uncertainty;
\item \kw{test.pot}, a global hard core and an initial interaction for each
  pair, $\epsilon = -e\,T\,[1.5\ln\ell + \ln(p/(1-p))]$ with $\ell$ the
  separation along the chain and $e$ the energy scale (the first term
  compensates the contact probability of an ideal chain);
\item \kw{test.par}, a parameter file with flip and pivot moves at $T=1$ and
  the optimization (\kw{op_minim sample}) over many runs.
\end{itemize}
Run
\begin{Verbatim}
  montegrappa test.pol test.pot test.par
\end{Verbatim}
After each run the interaction matrix is changed so that the contact
probabilities, reweighted from the frames of the run, approach the
experimental ones; the new matrix (\kw{newpot_<run>.pot}) is used in the
next run. \kw{restraints_<run>.dat} compares the target and the
back-calculated probability of each restraint (plot columns 2 and 3), and
\kw{chi2.dat} shows the progress of the fit. \kw{results/} contains the
input files and the potential after 15 iterations.

\subsection{Tutorial 3: parallel tempering and the specific heat of a hairpin}
\emph{Programs:} \kw{grappino}, \kw{montegrappa_mpi}, \kw{mhistogram}.
\emph{System:} residues 41--56 of protein G (\kw{hairpin.pdb}).

\kw{hairpin.in} builds a Go model of all heavy atoms: the N, C$_\alpha$, C
backbone and the sidechains as in the PDB (one rotamer each, no library),
with native dihedrals. \kw{hairpin.par} runs $10^7$ steps of parallel tempering with 8
replicas between $T=0.4$ and $1.8$, exchanging every 1000 steps:
\begin{Verbatim}
  grappino hairpin.in
  mpirun -np 8 montegrappa_mpi hairpin.pol hairpin.pot hairpin.par
\end{Verbatim}
Each replica writes \kw{energy_proc<N>.ene}, \kw{traj_proc<N>.pdb} and
\kw{last_proc<N>.pol}; the exchanges are in \kw{exchange_proc<N>.dat}.
After removing the first 10\% of each energy file as equilibration (e.g.\
\kw{sed -e '1,100d' energy_proc0.ene > energy_cut0.ene}),
\begin{Verbatim}
  mhistogram mhistogram.in
\end{Verbatim}
combines the energies of all temperatures with the multiple-histogram
method (\kw{nohisto}: the files are energy series, column 2) and writes in
\kw{results.dat} the average energy and the specific heat from $T=0.4$ in
100 steps of 0.014. The peak of the specific heat (at $T\simeq0.88$ in
\kw{results/}) marks the folding transition of the hairpin; \kw{dumb_e.dat} contains the plain averages at
each simulated temperature, to check the fit.

\subsection{Tutorial 4: optimization of a potential with parallel tempering}
\emph{Programs:} \kw{grappino}, \kw{montegrappa_mpi}.
\emph{System:} a helix of 15 alanines (\kw{polyala.pdb}).

\kw{polyala.in} builds a model with sidechains from the rotamer library
(\kw{rotamers}, \kw{rotfile ../../lib/rotamers.lib}, plus the rotamer of
the PDB with \kw{pdb_rot}), atom types from \kw{lib/atomtypes.1.lib}, the
peptide bond locked (\kw{locked_atoms N}), a weak Go potential and native
dihedrals. With \kw{op_file polyala.op} and \kw{op_kind GO_DIST_CA} it also
writes the restraints: the native distances between C$_\alpha$ atoms.
\kw{polyala.par} runs 5 iterations (\kw{nrun 5}) of $3\cdot10^6$ steps of
parallel tempering with 8 replicas between $T=1$ and $2.4$; after each one
the matrix is optimized at \kw{op_T 1.0} with frames recorded every 5000
steps after the first 200000. With \kw{always_restart} and
\kw{record_native} every iteration starts from, and includes, the native
structure.
\begin{Verbatim}
  grappino polyala.in
  mpirun -np 8 montegrappa_mpi polyala.pol polyala.pot polyala.par
\end{Verbatim}
The outputs carry the run and the replica in their name
(\kw{traj_run<M>_proc<N>.pdb}); \kw{chi2.dat} should decrease over the
iterations, and the RMSD of the replicas from the native helix shows how
the optimized potential stabilizes it (in \kw{results/}, from 0.5~\AA{} for
the coldest replica to 3~\AA{} for the hottest).

\subsection{Tutorial 5: adaptive simulated tempering of a random heteropolymer}
\emph{Programs:} \kw{montegrappa} built with \kw{version=STEMPERING}.
\emph{System:} a chain of 20 beads with a random matrix of interactions
(\kw{random_polymer.pot}).

The polymer and potential files are given. \kw{random_polymer.par} runs
$10^9$ steps of flip and pivot moves with adaptive simulated tempering:
starting from the single temperature $T=2$, the program collects energy
histograms, every $2\cdot10^7$ steps computes the density of states with
the multiple-histogram method (keeping all past histograms,
\kw{st_keepall}) and adds temperatures and adjusts their weights so that
the walk covers the range down to the target $T=0.6$.
\begin{Verbatim}
  montegrappa random_polymer.pol random_polymer.pot random_polymer.par
\end{Verbatim}
\kw{dumb.dat} and \kw{thermdodyn.dat} contain the average energy as a
function of temperature (columns 1:2), the latter also the specific heat
(columns 1:4); \kw{energy.jpeg} and \kw{Cv.jpeg} show the curves obtained
with this version: the ladder ends with five temperatures between 2 and
0.6, and the multiple-histogram fit passes through the average energy at
each of them.

\subsection{Tutorial 6: simulated tempering with optimization of the potential}
\emph{Programs:} \kw{montegrappa} built with \kw{version=STEMPERING}.
\emph{System:} a segment of 20 beads of DNA with 4 types of beads.

As tutorial 5, but the run also optimizes the interaction matrix of the 4
types so that beads 0 and 3 (the single restraint in \kw{dna.op}) are in
contact with probability 0.75 at $T=1.5$. The temperature ladder starts
from $T=5$ and adapts down to the target $T=1.5$; once it has reached it,
a frame is harvested at the target temperature every \kw{st_printpdb}
steps, and the optimization at the end of each of the two runs uses those
frames.
\begin{Verbatim}
  montegrappa dna.pol dna.pot dna.par
\end{Verbatim}
After the two runs, \kw{chi2.dat} shows the minimization of $\chi^2$.

%======================================================================
\section{Limits}
Compile-time limits are in \kw{include/montegrappa.h} (e.g.\ \kw{NRESMAX},
\kw{NATOMMAX}, \kw{NSIDEMAX}, \kw{NROTMAX}); the arrays are allocated with
the sizes found in the input, so the limits only bound the input checks.
\kw{grappino} reads at most 150 chains (\kw{NPOLMAX} in
\kw{include/grappino.h}). With MPI, the number of replicas is bounded by
\kw{NREPMAX}.

\end{document}
